\section{Matrimony}
\label{sec:matrimony}

Among the seven sacraments, Matrimony holds a singular place in salvation history: it alone did not arise after Adam's sin or await the Incarnation to begin its sacred calling. God inscribed marriage into the dawn of creation on the sixth day, forming humanity male and female in His own image and drawing the woman from the side of the sleeping Adam to stand beside him as bone of his bones and flesh of his flesh (Gen~1:27--28; 2:21--24). In that primeval sleep the Fathers beheld a prophetic icon of the Second Adam asleep in death upon the Cross, from whose pierced side flowed the water and blood that formed His Bride, the Church (Augustine, \emph{In Ioann. Evang.}~Tr.~120.2). What the Creator established in Paradise as a sacred covenant of nature, Jesus Christ freed from fallen hardness of heart and elevated into a true sacrament of the New Law---an indissoluble, grace-giving covenant by which one baptized man and one baptized woman image and share in the spousal union of Christ and His Church (Eph~5:25--32; Trent, Sess.~XXIV, can.~1; CIC~1983, c.~1055; CCEO~1990, c.~776).

\subsection{Scriptural Foundations: Creation, Covenant, Cana, and the Great Mystery}
\label{subsec:matrimony-scripture}

Sacred Scripture opens with a wedding and closes with a wedding. In Genesis~1, sexual difference and fruitfulness stand directly beside the divine image:
\begin{quotation}
And God created man to his own image: to the image of God he created him: male and female he created them. And God blessed them, saying: Increase and multiply, and fill the earth, and subdue it. (Gen~1:27--28)
\end{quotation}
Genesis~2 unfolds the personal communion and exclusive unity of that gift. Declaring ``It is not good for man to be alone: let us make him a help like unto himself'' (Gen~2:18), God builds Adam's rib into the woman and brings her to him, prompting Adam's greeting---``This now is bone of my bones, and flesh of my flesh'' (Gen~2:23)---and Scripture's permanent law of marriage:
\begin{quotation}
Wherefore a man shall leave father and mother, and shall cleave to his wife: and they shall be two in one flesh. (Gen~2:24)
\end{quotation}
Already before the Fall, four foundational truths shine together: humanity is ordered toward communion in two complementary sexes; husband and wife share equal personal dignity; their union is blessed with fruitfulness; and their covenant joins two alone into one indissoluble flesh.

Even while human sin obscured that primitive splendor through polygamy and the Mosaic tolerance of divorce, Old Testament revelation deepened the religious sanctity of marriage. On his wedding night, Tobias rises with Sara to pray:
\begin{quotation}
For we are the children of saints, and we must not be joined together like heathens that know not God\ldots\ Thou madest Adam of the slime of the earth, and gavest him Eve for a helper. And now, Lord, thou knowest, that not for fleshly lust do I take my sister to wife, but only for the love of posterity, in which thy name may be blessed for ever and ever. (Tob~8:5, 8--9)
\end{quotation}
The prophet Malachias condemns treachery against ``the wife of thy youth\ldots\ the wife of thy covenant'' because the Lord Himself stands as witness between them (Mal~2:14--16). Through Osee, Isaias, Jeremias, Ezechiel, and the Canticle of Canticles, God prepares humanity for the Christian mystery by revealing Himself as the faithful Bridegroom of Israel.

When the Word became flesh, He chose a wedding feast for His very first public sign. Invited with His Mother and disciples to Cana of Galilee, Jesus turned water for Jewish purification into abundant new wine: ``This beginning of miracles did Jesus in Cana of Galilee and manifested his glory. And his disciples believed in him'' (John~2:1--11). By His presence and first miracle, He ennobled wedlock and heralded the new wine of sacramental grace transforming natural love into supernatural charity (Leo~XIII, \emph{Arcanum Divinae}).

Tested later by Pharisees citing the Mosaic bill of divorce (Deut~24:1), Jesus reached behind Moses to Genesis and restored marriage's original unity and indissolubility by His supreme authority:
\begin{quotation}
Have ye not read, that he who made man from the beginning, made them male and female? And he said: For this cause shall a man leave father and mother, and shall cleave to his wife, and they two shall be in one flesh. Therefore now they are not two, but one flesh. What therefore God hath joined together, let no man put asunder\ldots\ Because Moses by reason of the hardness of your heart permitted you to put away your wives: but from the beginning it was not so. And I say to you, that whosoever shall put away his wife, except it be for fornication, and shall marry another, committeth adultery: and he that shall marry her that is put away, committeth adultery. (Matt~19:4--6, 8--9; cf.~Mark~10:6--12)
\end{quotation}
In Mark~10:11--12 and Luke~16:18, the Lord forbids remarriage after dismissal without qualification, obliging husband and wife equally. Hence patristic and conciliar teaching---stated plainly by St.~Bede (\emph{In Marci Evangelium}~10) and defined at Trent (Sess.~XXIV, can.~7)---reads the Matthean clause as allowing separation of bed and board in the case of adultery while never dissolving the bond or permitting another marriage during the spouse's lifetime.

St.~Paul repeats that exact command of the Lord: ``But to them that are married, not I, but the Lord, commandeth that the wife depart not from her husband. And if she depart, that she remain unmarried or be reconciled to her husband. And let not the husband put away his wife'' (1~Cor~7:10--11). Only death frees a spouse to marry again ``in the Lord'' (1~Cor~7:39; cf.~Heb~13:4). Speaking then by his own apostolic authority regarding two unbaptized spouses where one converts to Christ, Paul instructs the believer to stay if the unbeliever consents to live peacefully, yet adds: ``But if the unbeliever depart, let him depart. For a brother or sister is not under servitude in such cases. But God hath called us in peace'' (1~Cor~7:12--15)---the scriptural foundation of the Pauline privilege.

The summit of biblical teaching on marriage shines in Ephesians~5. Grounding domestic life in mutual reverence---``Being subject one to another, in the fear of Christ'' (Eph~5:21)---St.~Paul draws Christian marriage directly into Christ's Passion:
\begin{quotation}
Husbands, love your wives, as Christ also loved the church and delivered himself up for it: That he might sanctify it, cleansing it by the laver of water in the word of life: That he might present it to himself, a glorious church, not having spot or wrinkle or any such thing; but that it should be holy and without blemish\ldots\ For this cause shall a man leave his father and mother: and shall cleave to his wife. And they shall be two in one flesh. This is a great sacrament: but I speak in Christ and in the church. (Eph~5:25--27, 31--32)
\end{quotation}
Rendered \latin{Sacramentum hoc magnum est} in the Vulgate, Christian marriage is no mere poetic simile; it is caught up into Christ's sacrificial love and made its living sign. And because Christ's espousal to His Church finds consummation in heaven, every Christian wedding anticipates the Apocalypse's final chorus: ``For the marriage of the Lamb is come: and his wife hath prepared herself\ldots\ Blessed are they that are called to the marriage supper of the Lamb'' (Apoc~19:7, 9; cf.~21:2, 9).

\subsection{Patristic, Scholastic, and Magisterial Development}
\label{subsec:matrimony-tradition}

The Fathers of East and West defended marriage both against pagan licentiousness and against Gnostic or Manichaean contempt for bodily generation.

In the East, \textbf{St.~John Chrysostom} preached on Ephesians~5 with wonder at the hidden wisdom joining human love to Christ's Incarnation:
\begin{quotation}
For indeed, in very deed, a mystery it is, yea, a great mystery, that a man should leave him that gave him being, him that begat him, and that brought him up, and her that travailed with him and had sorrow\ldots\ and be joined to one who was never even seen by him and who has nothing in common with him, and should honor her before all others\ldots\ Such Moses prophetically showed it to be from the very first; such now also Paul proclaims it, where he saith, ``concerning Christ and the Church.'' (\emph{Homilies on Ephesians}~XX)
\end{quotation}
For Chrysostom, Christ left the Father's throne to descend to His Bride and become one spirit with her; harmony between husband and wife thus turns the Christian home into a little church and welds together families and cities alike.

In the West, \textbf{St.~Augustine} forged the enduring theological grammar of Christian marriage. In \emph{De civitate Dei}~XIV.21--23, he proves that bodily fecundity belongs to Eden rather than to the Fall:
\begin{quotation}
But that blessing upon marriage, which encouraged them to increase and multiply and replenish the earth, though it continued even after they had sinned, was yet given before they sinned, in order that the procreation of children might be recognised as part of the glory of marriage, and not of the punishment of sin. (\emph{De civitate Dei}~XIV.21)
\end{quotation}
Even after sin, marriage remains ``a kind of seed-bed of the city'' (\latin{quoddam seminarium est civitatis}): ``while the earthly city needs for its population only generation, the heavenly needs also regeneration to rid it of the taint of generation'' (\emph{De civitate Dei}~XV.16). Above all, in \emph{On the Good of Marriage} (\emph{De bono coniugali}) and in his literal commentary on Genesis (\emph{De Genesi ad litteram}~IX.7.12), Augustine formulated the classic \textbf{three goods of marriage} (\latin{bona matrimonii})---\textbf{fidelity} (\latin{fides}), \textbf{offspring} (\latin{proles}), and \textbf{sacrament} (\latin{sacramentum}). Fidelity excludes any union outside the conjugal bond; offspring means children lovingly welcomed, kindly nourished, and religiously educated; sacrament means indissolubility, so that spouses do not separate and neither may take another even for the sake of having children.

Medieval scholasticism---from Peter Lombard (\emph{Sententiae}~IV, dd.~26--42) to St.~Thomas Aquinas and the \emph{Supplementum} (qq.~41--68)---integrated Augustine's triad into a unified sacramental synthesis. Marriage is at once an office of nature, a civil-social bond, and, between baptized Christians, a grace-conferring sacrament of the New Law. Settling medieval debate over consent versus consummation, Pope Alexander~III, the scholastics, and the \textbf{Council of Florence} (\emph{Exsultate Deo}, 1439, \S16) established that the efficient cause of marriage is free mutual consent expressed in words of the present (\latin{per verba de praesenti}), while subsequent consummation renders the sacramental bond absolutely indissoluble.

Against sixteenth-century rejections of marriage's sacramentality, indissolubility, and ecclesiastical jurisdiction, the \textbf{Council of Trent} (Session~XXIV, 11~November 1563) defined Catholic dogma: Matrimony is truly one of the seven sacraments instituted by Christ and confers grace (can.~1); polygamy is forbidden (can.~2), and neither heresy, desertion, nor adultery dissolves the marriage bond (cans.~5, 7); the Church possesses native authority to establish diriment impediments, dispense from ecclesiastical impediments, and judge matrimonial causes (cans.~3--4, 12); sacred orders and solemn religious vows invalidate attempted marriage (can.~9); and consecrated virginity or celibacy stands objectively higher than the married state (can.~10). In the decree \emph{Tametsi} (Sess.~XXIV, \emph{Decretum de reformatione matrimonii}, ch.~1), Trent required canonical form before the parish priest and two or three witnesses to end clandestine marriages; and in ch.~10 of the same reform decree, governing solemnization during closed liturgical seasons, the Council admonished pastors that marriage is a holy reality to be treated in a holy manner (\latin{Sancta enim res est matrimonium et sancte tractandum}). The \emph{Roman Catechism} (1566, Part~II, \emph{On Matrimony}) drew the consequence cleanly: because Christian marriage signifies Christ's own unbreakable union with His Church, the marital bond can never be severed while both spouses live.

When nineteenth-century regalists and secular states attempted to sever Christian marriage into a state-governed ``civil contract'' plus an optional ecclesiastical ``sacrament,'' \textbf{Pope Leo~XIII} responded in his encyclical \emph{Arcanum Divinae} (10~February 1880), establishing the permanent principle of the \textbf{inseparability of valid contract and sacrament} between baptized Christians:
\begin{quotation}
A distinction, or rather severance, of this kind cannot be approved; for certain it is that in Christian marriage the contract is inseparable from the sacrament, and that, for this reason, the contract cannot be true and legitimate without being a sacrament as well. For Christ our Lord added to marriage the dignity of a sacrament; but marriage is the contract itself, whenever that contract is lawfully concluded\ldots\ Hence it is clear that among Christians every true marriage is, in itself and by itself, a sacrament; and that nothing can be further from the truth than to say that the sacrament is a certain added ornament, or outward endowment, which can be separated and torn away from the contract at the caprice of man. (\emph{Arcanum Divinae})
\end{quotation}
Whenever a valid matrimonial covenant arises between two baptized persons, it is by that very fact a sacrament of Christ (CIC~1983, c.~1055~\S2; CCEO~1990, c.~776~\S2).

The modern Magisterium unfolded this doctrine further. Pope Pius~XI's encyclical \emph{Casti Connubii} (31~December 1930) systematically expounded Augustine's three goods, defended unborn life and openness to procreation against artificial contraception and eugenics, and placed conjugal charity at the heart of marriage. The \textbf{Second Vatican Council} (\emph{Gaudium et Spes}~47--52; \emph{Lumen Gentium}~11) enriched the legal concept of contract with the biblical reality of personal \textbf{covenant} (\latin{foedus}): an intimate partnership of life and love rooted in irrevocable consent, in which conjugal love is taken up into divine charity, children are honored as marriage's supreme gift, and the Christian family lives as the \textbf{domestic church} (\latin{ecclesia domestica}). St.~John~Paul~II's Apostolic Exhortation \emph{Familiaris Consortio} (1981) developed this mission of the family as a community of life and love sharing in Christ's prophetic, priestly, and royal offices. Recent pontifical acts sharpened canonical clarity across both Codes: Benedict~XVI's \emph{Omnium in Mentem} (2009) removed the ``formal act of defection'' exception from CIC~1983 canons~1086~\S1, 1117, and~1124; Pope Francis's \emph{Mitis Iudex Dominus Iesus} and \emph{Mitis et misericors Iesus} (15~August 2015) reformed nullity tribunal procedure without touching substantive grounds of validity (such as c.~1095); and \emph{De concordia inter Codices} (31~May 2016) added CIC~1983 canon~1108~\S3 so that only a priest validly assists whenever an Eastern Catholic party marries.

\subsection{Sacramental Structure: Consent, East--West Discipline, and Indissolubility}
\label{subsec:matrimony-structure}

Because Matrimony consecrates a natural human covenant rather than employing an external substance such as water or oil, its sacramental structure is uniquely personal.

\textbf{Analogical Matter and Form.} Scholastic theology applies matter and form to Matrimony analogically: the spouses' mutual self-giving and acceptance of one another in their whole conjugal partnership constitute the quasi-matter, and the \textbf{legitimately manifested exchange of present marital consent} constitutes the form---joined in Eastern Catholic law to the sacred blessing of the priest (CCEO~1990, c.~828). Consent is an irreplaceable act of the will that no human power can supply (CIC~1983, c.~1057; CCEO~1990, c.~817). Grave incapacity of discretion (CIC~1983, c.~1095), simulation excluding marriage or an essential property such as fidelity, indissolubility, or openness to children (c.~1101~\S2), or external force and grave fear (c.~1103, authentically interpreted on 23~April~1987, AAS~79, p.~1132, as reaching even marriages of the unbaptized) prevents valid consent from coming into being.

\textbf{Subject.} The capable subject of sacramental Matrimony is \textbf{one eligible validly baptized man and one eligible validly baptized woman}, free of all diriment impediments of natural, divine, or ecclesiastical law (CIC~1983, cc.~1055, 1073--1094; CCEO~1990, cc.~776, 790--812). Impediments of natural or divine law---antecedent and perpetual impotence (CIC~1983, c.~1084~\S1, distinct from sterility, which neither prohibits nor nullifies, c.~1084~\S3), a prior existing marriage bond (c.~1085), or direct-line and second-degree collateral consanguinity (c.~1078~\S3)---cannot be dispensed by any human authority or sanated while the impediment lasts (c.~1163~\S2). Ecclesiastical impediments may be dispensed by competent authority, with sacred orders, public perpetual vow of chastity in a pontifical institute, and conjugicide reserved to the Apostolic See (CIC~1983, cc.~1078, 1087--1088; CCEO~1990, cc.~795, 804). Both spouses share equal duties and rights in conjugal life (CIC~1983, c.~1135). When one or both spouses are unbaptized, a valid union is an honorable \textbf{natural marriage}, intrinsically monogamous and lifelong by natural law, which becomes a sacrament \latin{ipso facto} if both spouses are subsequently baptized.

\textbf{Minister in Latin and Eastern Traditions.} East and West preserve two complementary expressions of one Catholic faith (CCC~1623):
\begin{itemize}
  \item In the \textbf{Latin Church}, the \textbf{spouses themselves} mutually confer the sacrament by expressing their free consent before the Church; the bishop, priest, or deacon (or, where authorized under CIC~1983, c.~1112, a delegated lay person) assists as the Church's \textbf{qualified official witness}, asking for and receiving their consent in the Church's name and imparting the nuptial blessing (CIC~1983, cc.~1108--1112).
  \item In the \textbf{Eastern Catholic Churches}, the \textbf{priest or bishop} receives the spouses' mutual consent and confers the sacrament through the sacred blessing---in Byzantine and sister traditions, the solemn \textbf{Crowning} of bridegroom and bride as king and queen of a new domestic church and martyrs of self-giving love (CCEO~1990, c.~828~\S1; CCC~1623). Because a priestly blessing is ordinarily required for validity, a deacon cannot validly assist at an Eastern marriage, a rule now mirrored in Latin law whenever an Eastern party marries (CIC~1983, c.~1108~\S3).
\end{itemize}

\textbf{Ordinary and Extraordinary Canonical Form.} Catholics are ordinarily bound for validity to \textbf{canonical form} before a competent sacred minister and two witnesses (CIC~1983, cc.~1108, 1117; CCEO~1990, cc.~828, 834); outside urgent danger of death, a diocesan bishop cannot dispense two Catholics from canonical form (Authentic Interpretation of CIC~c.~87~\S1, 5~July~1985, AAS~77, p.~771). In \textbf{danger of death}, or when a competent minister cannot be reached without grave inconvenience for at least a month, both Codes provide for \textbf{extraordinary canonical form} before two witnesses alone (CIC~1983, c.~1116; CCEO~1990, c.~832), with Eastern law directing that a priestly blessing still be sought as soon as possible afterward.

\begin{threestable}{Canonical Dimension}{Latin Church (CIC~1983)}{Eastern Catholic Churches (CCEO~1990)}
Minister of the sacrament & Spouses mutually confer the sacrament; bishop, priest, or deacon (or delegated lay witness, c.~1112) assists as official witness (CCC~1623) & Priest or bishop receives consent and confers the sacrament through the sacred blessing / Crowning (CCEO~c.~828~\S1; CCC~1623)\\
Essential canonical form & Consent exchanged before authorized ordinary, pastor, or delegated priest/deacon and two witnesses (cc.~1108, 1117) & Consent exchanged before local Hierarch, pastor, or delegated priest who imparts the sacred blessing, with two witnesses (c.~828); deacon cannot validly assist\\
Inter-ritual marriage & Under \emph{De concordia inter Codices} (2016), only a priest validly assists when an Eastern party marries (CIC~c.~1108~\S3; c.~1127~\S1) & Priestly blessing required for validity (CCEO~cc.~828, 834~\S2), harmonizing East and West across rites\\
Extraordinary form & Valid before two witnesses alone in danger of death or foreseen one-month unavailability of competent minister (c.~1116) & Valid before two witnesses alone under the same conditions (c.~832); priestly blessing to be sought as soon as possible\\
\end{threestable}

\textbf{Absolute Indissolubility, Nullity, and Bounded Dissolutions.} Catholic doctrine keeps three juridical realities strictly distinct:
\begin{enumerate}
  \item \textbf{Ratified and consummated marriage (\latin{ratum et consummatum}):} A valid sacramental marriage between two baptized persons consummated by a conjugal act apt for generation \textbf{can never be dissolved by any human power or for any cause other than death} (CIC~1983, c.~1141; CCEO~1990, c.~853). Not even the Pope can dissolve it.
  \item \textbf{Declaration of nullity and separation with bond remaining:} A declaration of nullity does not dissolve a valid marriage; it establishes that no valid bond ever arose from the beginning owing to an impediment, defect of consent, or lack of canonical form. For grave causes such as adultery or serious danger to soul or body, the Church permits physical \textbf{separation while the bond remains} (CIC~1983, cc.~1151--1155; CCEO~1990, cc.~863--866), without freedom to remarry.
  \item \textbf{Strictly bounded dissolutions of non-consummated or non-sacramental bonds:} Leaving \latin{ratum et consummatum} indissolubility inviolate, canon law recognizes three narrow dissolutions: (a) a \textbf{ratified but non-consummated marriage} (\latin{ratum sed non consummatum}) may be dissolved by the Roman Pontiff alone for a just cause (CIC~1983, c.~1142; CCEO~1990, c.~862); (b) under the \textbf{Pauline privilege} (1~Cor~7:12--15; CIC~1983, cc.~1143--1147; CCEO~1990, cc.~854--858), a natural marriage of two unbaptized persons dissolves in favor of the faith when one is baptized, the unbaptized party departs, and the baptized spouse contracts a new marriage; and (c) under the \textbf{Petrine privilege} (privilege of the faith), the Roman Pontiff may dissolve a non-sacramental marriage (at least one spouse unbaptized throughout) in favor of the faith and salvation of souls.
\end{enumerate}

\subsection{Sacramental Reality, Relational Bond, Grace, and Order of Ends}
\label{subsec:matrimony-effects}

\textbf{\latin{Res et sacramentum} and Permanent Relational Change (No Character).} Matrimony imprints \textbf{no indelible sacramental character}; death ends the earthly bond and leaves the surviving spouse free to marry again in the Lord (1~Cor~7:39). Its immediate, objective spiritual reality (\latin{res et sacramentum}) is the \textbf{perpetual and exclusive marital bond (\latin{vinculum})} joined by God Himself---a real, \textbf{permanent relational change} uniting husband and wife as two in one flesh (CCC~1638--1640; \emph{Roman Catechism}, Part~II, On Matrimony). During their shared earthly life, this bond endures independently of shifting emotions or subsequent faults as the living icon of Christ's irrevocable faithfulness to His Church.

\textbf{\latin{Res tantum} and Sacramental Grace.} In spouses free of mortal sin (and reviving upon repentance if an obstacle initially blocked fruitfulness), Matrimony confers an increase of sanctifying grace joined to the \textbf{proper sacramental grace of the married state} (Trent, Sess.~XXIV, can.~1; \emph{Gaudium et Spes}~48; CCC~1641--1642). This grace heals and elevates natural conjugal love into supernatural charity, strengthens lifelong fidelity through trial and age, aids parents in welcoming and forming children in the faith, and consecrates daily family life as a \textbf{domestic church} (\emph{Lumen Gentium}~11). As Leo~XIII wrote (\emph{Arcanum Divinae}), the duties of marriage are neither few nor light, yet by the strength of this sacrament they become not merely bearable but joyful.

\textbf{Order of Ends: Classical Hierarchy and Conciliar Co-Ordering.} Catholic theology harmonizes classical scholastic terminology with the personalist deepening of Vatican~II and the revised Codes. Following St.~Augustine and St.~Thomas (\emph{Supplementum}, q.~41, a.~1), Pius~XI (\emph{Casti Connubii}) and the 1917 Code (c.~1013~\S1) termed the \textbf{procreation and education of offspring} marriage's primary end, and \textbf{mutual help} (\latin{mutuum adiutorium}) with the \textbf{remedy of concupiscence} (\latin{remedium concupiscentiae}) its secondary ends. Vatican~II (\emph{Gaudium et Spes}~48--50) and current universal law (CIC~1983, c.~1055~\S1; CCEO~1990, c.~776~\S1) teach that the matrimonial covenant is by nature \textbf{ordered to the good of the spouses (\latin{bonum coniugum}) and to the procreation and education of offspring (\latin{prolis generatio et educatio})}. Rather than setting spousal love and fecundity against one another, this co-ordered expression discloses their organic unity: authentic conjugal self-gift is inherently life-giving, and new life finds its proper sanctuary within the faithful communion of father and mother. Hence even when physical sterility prevents conception without fault, a valid marriage preserves its full dignity, indissolubility, and sanctifying grace (\emph{Gaudium et Spes}~50; CIC~1983, c.~1084~\S3).

Across all four teleological dimensions, the ends of Matrimony stand in clear order:
\begin{itemize}
  \item \textbf{Primary proper ends:} the good of the spouses in faithful, lifelong communion of persons and the procreation and Christian education of offspring (CIC~1983, c.~1055~\S1; CCEO~1990, c.~776~\S1; \emph{Gaudium et Spes}~48--50).
  \item \textbf{Intrinsic secondary ends:} mutual aid in spiritual and temporal life, the right ordering and healing of concupiscence (\emph{ST}~III, q.~65, a.~1), the sanctification of the home as a domestic church (\emph{Lumen Gentium}~11), and the stability of human society (\emph{De civitate Dei}~XV.16; \emph{Arcanum Divinae}).
  \item \textbf{Contingent ends:} concrete works of family hospitality, Christian stewardship, and civic service shaped by each household's vocation and era.
  \item \textbf{Ultimate end:} eternal union of spouses and children with the Triune God at the Marriage Supper of the Lamb (Apoc~19:7--9; 21:2), where the earthly sign gives way to the eternal reality it signified.
\end{itemize}

\begin{threestable}{Structural Dimension}{Doctrinal Substance}{Canonical and Liturgical Specification}
Analogical matter & Mutual self-giving and acceptance of one eligible baptized man and one eligible baptized woman & No physical element is used; hylomorphic categories apply analogically to the personal conjugal donation\\
Analogical form & Legitimately manifested exchange of present marital consent & Irreplaceable act of will (CIC~c.~1057; CCEO~c.~817); Eastern Catholic law ordinarily requires priestly blessing for validity (CCEO~c.~828)\\
Minister & Latin: spouses mutually confer the sacrament; East: priest or bishop confers the sacred blessing & Latin sacred minister assists as qualified witness (CIC~c.~1108); Eastern priest confers the Crowning/blessing (CCC~1623; CIC~c.~1108~\S3 for Eastern parties)\\
Subject & One validly baptized man and one validly baptized woman, free of diriment impediments & Natural-law impediments (impotence, prior bond, direct-line consanguinity) cannot be dispensed; sterility does not invalidate (CIC~c.~1084~\S3)\\
\latin{Res et sacramentum} & Perpetual, exclusive marital bond (\latin{vinculum}), image of Christ and the Church & Permanent relational change; \textbf{no sacramental character}; ratified and consummated bond dissolves only by death (CIC~c.~1141; CCEO~c.~853)\\
\latin{Res tantum} & Increase of sanctifying grace plus spousal grace perfecting love, fidelity, and parenthood & Builds up the home as \latin{ecclesia domestica} (\emph{LG}~11; \emph{GS}~48); valid marriage of two baptized persons is inseparably a sacrament (\emph{Arcanum Divinae}; CIC~c.~1055~\S2)\\
\end{threestable}
